\chapter{Conclusion}

The website is the web face of an open-source health platform that had only a mobile app until now. It puts everything the site offers in one place in a browser, with nothing to install.

A visitor can read health articles from MedlinePlus, the NHS, the World Health Organization and its India office, the UK Health Security Agency, and India's national centre for vector-borne diseases, with the publisher and licence shown on each one. They can ask the health assistant a question and get an answer that comes only from those articles, with the sources listed underneath; when the site has nothing on the question, it says so instead of guessing. They can describe symptoms and be told what to do next, with urgent instructions if the description includes warning signs, and with questions rather than a diagnosis. They can read a chikungunya page written from those named sources, open health videos on the official channels those authorities publish on, and read outside questions and answers clearly marked as public Q&A alongside links to established communities elsewhere.

With an account, a person can record their own readings and see them over time, set their own goals, get suggestions based on the interests they chose, save articles and episodes, and download or delete everything the site holds about them. The community area lets a person post about a disease, reply, react and follow, and there is a moderation page for the operator. The site does not invent posts or testimonials.

The site is honest about what it is for. It does not diagnose, give medication advice, or estimate risk, and no doctor has reviewed its information; the privacy and terms pages say what is actually stored and for how long. The answers, the emergency warnings, signing in, keeping one person's information away from another's, and the community were all checked, and those checks stay with the site so they can be run again. The assistant will not compare chikungunya and dengue from its sources, so that comparison is made on the chikungunya page instead.

Taken together, the website shows that a health education site can be useful without pretending to be a clinician, as long as it shows where its answers come from and says plainly when it does not know. It is a working base a maintainer can keep building on: grow the library of articles, cover the comparisons it still leaves to a page rather than an answer, and record a real doctor's review in the place the site already provides for that, once one is available.

Possible next steps, in plain order:

\begin{itemize}
  \item Bring in more topics and more languages, and decide on a regular check so no article quietly drifts out of date.
  \item Fix the questions the assistant still leaves unanswered, especially the ones where it is comparing two illnesses and one of them is missing from its source list.
  \item Make the source checking clearer over time, so each answer can be followed back to what it was made from and any source changes are not hidden.
  \item Put personal records on a proper shared store with a real account life, so a person can change, download, remove or recover their data instead of it living only in one browser.
  \item Add symptom guidance for more problems over time, while keeping the same shape: act on urgent signs first, then a care level, then the sources.
  \item Write down clearly which parts would need a clinical review before they could be used in any real care setting, so health education stays separate from medical advice.
\end{itemize}
